% 9 câu đã gỡ khỏi ngân hàng (chuyen_dang) — chép lại vào file dạng để khôi phục

% từ: cau-hoi/lop10/0D3_Ham_so_bac_hai_va_do_thi/0D3-2-4_Bai_toan_ve_su_tuong_giao.tex
\begin{bt}%[0D3B2-4]%[ma:79c02afe]
	Giải  các phương trình sau
	\begin{enumerate}
		\item $\sqrt{x^2  + 3x - 3}=2x - 3$.
		\item $\sqrt{3x^2 - 2x + 3}=\sqrt{x^2  + 3x + 1}$.
	\end{enumerate}
	 \loigiai{
		\begin{enumerate}
			 \allowdisplaybreaks{
			\item \begin{eqnarray*}
			& &\sqrt{x^2 + 3x - 3}=2x -  3 \\
			 &\Leftrightarrow &\heva{& 2x-3\geq0 \\ & x^2+3x-3=(2x-3)^2}\\
			&\Leftrightarrow&\heva{& 2x\geq 3 \\ & x^2+3x-3=4x^2-12x+9}\\ 
			 &\Leftrightarrow &\heva{& x\geq\dfrac{3}{2} \\ & 3x^2-15x+12=0}\\
			&\Leftrightarrow&\heva{& x\geq\dfrac{3}{2} \\ & \hoac{& x=1  \mbox{ (loại)} \\ & x=4 \mbox{ (nhận)}.}}
		\end{eqnarray*}
		Vậy $x=4$. 
		\item \begin{eqnarray*}
		& &\sqrt{3x^2 - 2x + 3}=\sqrt{x^2 + 3x +  1}\\
		&\Leftrightarrow &\heva{&  x^2 + 3x + 1\geq0 \\ & 3x^2-2x+3=x^2 + 3x + 1}\\
		&\Leftrightarrow&\heva{& \hoac{& x\geq\dfrac{-3+\sqrt{5}}{2} \\ & x\leq\dfrac{-3-\sqrt{5}}{2}} \\  & 2x^2-5x+2=0}\\
		&\Leftrightarrow&\heva{& \hoac{& x\geq\dfrac{-3+\sqrt{5}}{2} \\ & x\leq\dfrac{-3-\sqrt{5}}{2}} \\  & \hoac{& x=2 \\ & x=\dfrac{1}{2}} \mbox{ (nhận)}.}
	\end{eqnarray*}}
	Vậy $x=2$,  $x=\dfrac{1}{2}$.
	\end{enumerate}
} 
\end{bt}

% từ: cau-hoi/lop11/1D1_Ham_so_luong_giac_va_phuong_trinh_luong/1D1-2-2_Xet_dau_gia_tri_luong_giac_Tinh_gia_tri_luong.tex
\begin{ex}%[1D1N2-2]%[ma:4ae6bf6d]
	Cho  cấp số cộng $(u_n)$ có $u_1=-3$, $u_7=21$. Khi đó công sai $d$ là
	\choice 
		 {$3$}
		{\True  $4$}
		 {$24$}
		{$\dfrac{24}{7}$} 
	\loigiai{Áp dụng công  thức $u_n=u_1+(n-1)d$. Khi đó
		 $$
		u_7=u_1+6d \Rightarrow d=\dfrac{u_7-u_1}{6}=\dfrac{21-(-3)}{6}=4. 
	 $$}
\end{ex}

% từ: cau-hoi/lop11/1D1_Ham_so_luong_giac_va_phuong_trinh_luong/1D1-4-1_Cau_hoi_ly_thuyet.tex
\begin{ex}%[1D1V4-1]%[ma:310fb045]
 Cho dãy số $(u_n)$ với $u_n = -3n  + 6$.
  \choiceTF
  {Trong dãy  số $(u_n)$ thì $n$ nhận giá trị là các số nguyên}
  {\True Số hạng  thứ $5$ của dãy là $-9$}
   {Dãy số $(u_n)$ là dãy số bị chặn}
  {\True Dãy số $(u_n)$ là  dãy số giảm}
 \loigiai{ 
  Ta có $u_{n+1} =  -3(n+1) + 6$ và $u_n = -3n + 6$.\\
  Suy ra  $k = u_{n+1} - u_n = -3$.\\
  Do đó dãy số $(u_n)$ giảm. \\ 
  Mặt khác $u_1 = 3$ nên dãy  số bị chặn trên bởi $3$ và dãy số không bị chặn dưới suy ra dãy số không bị chặn.
  \begin{itemchoice}
   \itemch \textbf{Sai.} $n\in\mathbb{N}^*$. 
   \itemch  \textbf{Đúng.} $u_5=-3\cdot 5+6=-9$.
   \itemch \textbf{Sai.} Dãy số $(u_n)$ không bị chặn. 
   \itemch \textbf{Đúng.} Dãy số $(u_n)$ là dãy số  giảm.
 \end{itemchoice}}
\end{ex}

% từ: cau-hoi/lop11/1D3_Gioi_han_Ham_so_lien_tuc/1D3-2-1_Cau_hoi_ly_thuyet.tex
\begin{ex}%[1D3K2-1]%[ma:66f1dc28]
	Cho dãy số $(u_n)$ với $\heva{&u_1=5 \\& u_{n+1}=u_n+n}$. Số hạng tổng quát của  dãy số này là
	\choice 
		 {$u_n=\dfrac{(n-1)n}{2}$}
		 {\True $u_n=5+\dfrac{(n-1)n}{2}$}
		{$u_n=5+\dfrac{(n+1)n}{2}$} 
		{$u_n=5+\dfrac{(n+1)(n+2)}{2}$} 
	\loigiai{Từ tính chất  của dãy ta có
		\begin{align*}
			&u_1=5 \\ 
			&u_2=u_1+1\\ 
			&u_3=u_2+2\\ 
			 &\cdots=\cdots\\
			&u_n=u_{n-1}+(n-1) 
		\end{align*}
		 Suy ra
		\begin{align*}
			&  u_1+u_2+\cdots +u_n=5+(u_1+u_2+\cdots +u_{n-1})+[1+2+\cdots +(n-1)]\\
			\Leftrightarrow \; & u_n=5+\dfrac{(n-1)n}{2}. 
		\end{align*}
	 }
\end{ex}

% từ: cau-hoi/lop11/1D3_Gioi_han_Ham_so_lien_tuc/1D3-2-1_Cau_hoi_ly_thuyet.tex
\begin{ex}%[1D3Y2-1]%[ma:ce64d18e]
	Cho dãy số $(u_n)$ với các số hạng đầu là $5;10;15;20;25;\ldots$. Tìm số hạng tổng quát của dãy số này. 
	\loigiai{Từ tính chất các số hạng của dãy suy ra số hạng tổng quát của dãy $(u_n)$ là $u_n=5n$.} 
\end{ex}

% từ: cau-hoi/lop11/1D3_Gioi_han_Ham_so_lien_tuc/1D3-2-3_PP_dat_thua_so_chung_ket_qua_huu_han.tex
\begin{ex}%[1D3B2-3]%[ma:27e13b8f]
	Trong các dãy số $( u_n)$ cho bởi số  hạng tổng quát $u_n$ sau, dãy số nào là dãy số giảm?
	 \choice
		{  \True $u_n=\dfrac{1}{{{2}^{n}}}$}
		{ $u_n=\dfrac{3n-1}{n+1}$} 
		{ $u_n={{n}^2}$} 
		 { $u_n=\sqrt{n+2}$}
	 \loigiai{
	Ta có $u_n=\dfrac{1}{{{2}^{n}}}$$<\dfrac{1}{{{2}^{n+1}}}={{u}_{n+1}}$  $\forall n\in {{\mathbb{N}}^{*}}$}
\end{ex}

% từ: cau-hoi/lop11/1D3_Gioi_han_Ham_so_lien_tuc/1D3-2-3_PP_dat_thua_so_chung_ket_qua_huu_han.tex
\begin{ex}%[1D3Y2-3]%[ma:f1cb87e9]
	Cho dãy số $( u_n)$ với $u_n={{n}^2}+1$. Xét tính tăng giảm  của dãy số.
	\loigiai{ 
	Ta có  $u_{n+1}-u_{n}=\left(n+1\right)^2+1-n^2-1=2n+1>0, \forall n\in {{\mathbb{N}}^{*}}$.\\ Vậy $u_n$ là dãy tăng.}
\end{ex}

% từ: cau-hoi/lop11/1D3_Gioi_han_Ham_so_lien_tuc/1D3-3-2_Tinh_lien_tuc_the_hien_qua_do_thi.tex
\begin{ex}%[1D3B3-2]%[ma:999cf900]
	 Cho cấp số cộng $({{u}_{n}})$ có ${{u}_{3}}=15$ và $d=-2$. Tìm ${{u}_{n}}$
	 \choice
		{ ${{u}_{n}}=\dfrac{3}{2}{{n}^2}-4$} 
		{ ${{u}_{n}}=-3n-17$} 
		{ \True  ${{u}_{n}}=-2n+21$}
		{  ${{u}_{n}}=-\dfrac{3}{2}n+12$}
	\loigiai{ 
		Ta có $\heva{ 
		&  15={{u}_{3}}={{u}_{1}}+2d \\
		 & d=-2 \\
		}\Leftrightarrow  \heva{
		 & {{u}_{1}}=19 \\
		 & d=-2 \\
	}\Rightarrow  {{u}_{n}}={{u}_{1}}+(n-1)d=-2n+21$.}
\end{ex}

% từ: cau-hoi/lop11/1D3_Gioi_han_Ham_so_lien_tuc/1D3-3-5_Bai_toan_phuong_trinh_co_nghiem.tex
\begin{ex}%[1D3K3-5]%[ma:9785bcca]
	Tổng $n$ số hạng đầu tiên của một cấp số cộng  là ${{S}_{n}}=\dfrac{3{{n}^2}-19n}{4}$ với $n\in {{\mathbb{N}}^{*}}$. Tìm số hạng đầu tiên ${{u}_{1}}$ và công sai $d$ của cấp số cộng đã cho.
	 \choice
		 { \True ${{u}_{1}}=-4;d=\dfrac{3}{2}$}
		{ ${{u}_{1}}=\dfrac{5}{2};d=\dfrac{1}{2}$} 
		{ ${{u}_{1}}=-\dfrac{3}{2};d=-2$} 
		{  ${{u}_{1}}=2;d=-\dfrac{1}{2}$}
	 \loigiai{
		Ta  có $$\dfrac{3{{n}^2}-19n}{4}=\dfrac{3}{4}{{n}^2}-\dfrac{19}{4}n={{S}_{n}}=n{{u}_{1}}+\dfrac{{{n}^2}-n}{2}d=\dfrac{d}{2}{{n}^2}+\left({{u}_{1}}-\dfrac{d}{2}\right)n
		\Leftrightarrow \heva{ 
		& \dfrac{d}{2}=\dfrac{3}{4} \\ 
		 & {{u}_{1}}-\dfrac{d}{2}=-\dfrac{19}{4} \\
		}\Leftrightarrow  \heva{
		 & {{u}_{1}}=-4 \\
		&  d=\dfrac{3}{2}. \\
	 }$$}
\end{ex}
